% additional file for Chapter 1
\chapter{Introduction}
\section{Overview}
The unprecedented surge in global data traffic, driven largely by the widespread deployment of 5G cellular networks and Fiber-to-the-Home (FTTH) access architectures, has fundamentally reshaped the telecommunications landscape. To satisfy this monumental and continuous demand without incurring the prohibitive costs of laying new physical fiber, network operators are compelled to maximize the transmission capacity of existing optical resources. This is primarily achieved by aggregating multiple communication channels into a single optical fiber using Wavelength Division Multiplexing (WDM) and, more specifically, Dense Wavelength Division Multiplexing (DWDM), which increases fiber efficiency by transmitting independent signals at different wavelengths.
\begin{figure}[htbp]
    \centering
    \includegraphics[width=0.6\textwidth]{internet_user_report_oct_25_datareportal.png}
    \caption{Number of internet users over the time}
    \label{fig:internet_user_report_oct_25_datareportal}
\end{figure}
However, as individual fiber capacity nears its theoretical Shannon limits, spectral multiplexing alone is often insufficient to satisfy massive localized traffic requirements. Consequently, network operators increasingly deploy "multifiber" or parallel-link configurations—utilizing multiple physical fiber pairs between adjacent nodes—to achieve a multiplicative increase in total throughput. While protocols like Link Aggregation (LAG) are traditionally utilized at higher network layers to duplex or group these connections for redundancy and increased bandwidth, managing these resources at the optical layer introduces significant complexity.

To handle the management burden of these high-capacity, multifiber infrastructures, the optical control plane has transitioned from static, manual configurations to highly automated, centralized management systems empowered by Software-Defined Networking (SDN). As DWDM technologies have matured, the traditional fixed-grid spectrum allocation—which assigns rigid, uniform channel spacing—has become a structural bottleneck. To overcome this, flexible-grid (flex-grid) architectures have emerged as the industry standard. By adhering to ITU-T G.694.1 standards, which utilize a precise 6.25 GHz slot granularity with 12.5 GHz slot width, flex-grid networks allow for elastic, demand-driven bandwidth allocation that vastly improves overall spectral efficiency \cite{itu694}.
\begin{figure}[htbp]
    \centering
    \includegraphics[width=0.7\textwidth]{example_of_flex_grid.png}
    \caption{An example of use of flexible the grid}
    \label{fig:example_usage_of_flex_grid}
\end{figure}
The deployment of flex-grid systems in long-haul backbone communications introduces substantial algorithmic complexity. In multi-hop networks, an optical signal must maintain both spectrum continuity—utilizing the same frequency across all links—and spectrum contiguity—utilizing adjacent frequency slots—from source to destination. This Routing and Spectrum Allocation (RSA) challenge is critically exacerbated in topologies featuring multiple parallel optical links between network elements, such as adjacent Reconfigurable Optical Add-Drop Multiplexers (ROADMs). When parallel links are introduced to a multi-hop architecture, the number of viable routing permutations increases exponentially. Consequently, traditional SDN controllers struggle to efficiently evaluate the availability of every 6.25 GHz slot across all possible path combinations, leading to severe computational bottlenecks and delayed provisioning times.

To resolve this critical scalability issue, this thesis proposes a novel approach by introducing the Parallel Optical Controller, a stateful microservice integrated directly into the TeraFlowSDN ecosystem. By leveraging advanced graph modeling techniques and a highly optimized bitwise spectrum intersection algorithm, this proposed method efficiently isolates available spectrum blocks across complex parallel routes. Ultimately, this solution effectively neutralizes the parallel link bottleneck, ensuring rapid, reliable, and highly scalable lightpath provisioning for next-generation optical networks.
\section{Motivation}
While extensive research has been conducted on Routing and Spectrum Allocation (RSA) in Elastic Optical Networks (EONs)\cite{algorithmsRSATaxonomy}\cite{efficientDynamicRSAEon},particularly concerning complex multi-hop lightpath provisioning\cite{heuristicWSAMultihop}—there remains a notable gap in the literature regarding the efficient management of multiple parallel optical links between adjacent devices. In traditional centralized control planes, such as those managed by other SDNs, parallel links are typically treated as isolated, individual entities rather than aggregated resources\cite{novelRoutingWSA}. Consequently, when a lightpath request is processed, the RSA algorithm is forced to evaluate the spectrum continuity of each parallel link independently\cite{efficientDynamicRSAEon}. This legacy approach is highly inefficient, as it triggers redundant, computationally heavy, back-and-forth calculations that severely degrade the controller's performance and scalability in dense optical topologies.

The primary motivation of this thesis is to fundamentally resolve this parallel link computational bottleneck within the emerging TeraFlowSDN ecosystem \cite{aboutTFS}. By addressing relative algorithmic inefficiency, the research aims to design a system capable of evaluating multiple parallel paths concurrently, significantly reducing the computational overhead required to find available flex-grid spectrum.

Furthermore, theoretical RSA optimization cannot be solved as just a theoretical math problem or a simulation using static data; its efficacy is strictly dependent on the acquisition of precise, real-time telemetry from the underlying physical layer. While modern optical networks are rapidly transitioning toward vendor-agnostic architectures unified by standardized data models such as OpenConfig\cite{openconfigYang} and OpenROADM\cite{openroadmYang}, a significant barrier remains in the practical deployment of centralized algorithms: hardware heterogeneity. Although these industry standards define common data models using YANG, individual equipment vendors frequently implement them with subtle variations in configuration hierarchies and telemetry reporting. Consequently, extracting uniform network state information—such as operational frequency ranges and transport capabilities—via a singular filtering procedure is often unfeasible in multi-vendor environments. This research addresses this disparity by bridging the gap between abstract path computation and physical hardware limitations through the extensive enhancement of device drivers within the TeraFlowSDN ecosystem. By refining the default driver logic and implementing a tiered approach that utilizes both standardized procedures and specialized, device-based handlers, this thesis establishes a more robust framework for multi-vendor support. These enhancements ensure that the proposed RSA algorithm is fed by accurate, model-driven telemetric data, effectively reconciling high-level computational logic with the practical complexities of diverse optical hardware management.
\section{Research Objectives}
To systematically address the challenges of Routing and Spectrum Allocation (RSA) over parallel optical links, this thesis is guided by the following core research objectives:
\begin{itemize}
    \item \textbf{Analysis of Flexible-Grid Standards:} To investigate the operational principles of flexible-grid optical networks, specifically adhering to the ITU-T G.694.1 standards regarding spectrum slot width and 6.25 GHz granularity.
    \item \textbf{Evaluation of SDN Optical Intent Provisioning:} To study how optical intents and lightpath requests are currently handled by popular, industry-standard Software-Defined Networking (SDN) controllers, identifying their limitations in dense, multi-link topologies.
    \item \textbf{Examination of TeraFlowSDN Device Drivers:} To conduct a comprehensive study of the telemetry and configuration mechanisms within TeraFlowSDN, with a specific focus on analyzing the OpenConfig and OpenROADM device drivers.
    \item  \textbf{Assessment of the Existing Optical Controller:} To critically evaluate the official OpticalController microservice currently operating within the TeraFlowSDN ecosystem, identifying its architectural constraints when processing complex optical intents.
    \item \textbf{Review of RSA Algorithmic Approaches:} To explore and analyze contemporary algorithmic strategies used to solve the RSA problem in multi-hop optical networks, evaluating their computational efficiency and spectrum assignment mechanisms.
    \item \textbf{Formulation of an Optimal Solution:} To synthesize the findings from the aforementioned studies and propose an optimal, highly scalable RSA framework that efficiently resolves the computational bottlenecks associated with parallel optical links.
\end{itemize}
\section{Thesis Contributions}
Building upon the established research objectives, this thesis makes several tangible contributions to the field of optical Software-Defined Networking (SDN), specifically addressing the computational bottlenecks associated with routing over parallel links. The core contributions of this work are summarized as follows:

\begin{itemize}
    \item \textbf{Enhancement of Optical Device Drivers:} A significant extension of the existing TeraFlowSDN device drivers, specifically the OCDriver (for OpenConfig devices) and the OpenROADM driver. This contribution includes the implementation of custom filters to extract vendor-agnostic telemetry—such as minimum and maximum operating frequencies, transport types (e.g., OCH, OMS), and flexible-grid parameters—directly from XML configurations, seamlessly integrating this data into an extended Context database schema.
    \item \textbf{Design of an Optimized RSA Algorithm:} The formulation and implementation of a highly scalable Routing and Spectrum Allocation (RSA) algorithm tailored for topologies with multiple parallel links. This includes utilizing the networkX library to construct multi-directed graphs that accurately represent parallel edges, alongside the introduction of a novel bitmap-based bitwise intersection method designed to rapidly compute the availability of contiguous 6.25 GHz spectrum slots across multiple hops.
    \item \textbf{Development of the ParallelOpticalController Microservice:} The creation, deployment, and integration of a novel, stateful microservice—the ParallelOpticalController—within the TeraFlowSDN microservices architecture. This dedicated controller autonomously processes complex optical intents, executing the advanced RSA algorithm to efficiently discover primary shortest paths and alternative route combinations without overwhelming the control plane.
    \item \textbf{Advanced Network Visualization and WebUI Integration:} To comprehensively upgrade the TeraFlowSDN Web User Interface (WebUI) to reflect the complex physical and logical states of the optical network. This objective focuses on developing visual representations of device endpoint capabilities, real-time spectrum slot availability (based on ITU-T 6.25 GHz granularity), multi-hop topological links, and the dynamic output of the RSA computation.
\end{itemize}
By fulfilling these objectives, this research aims to deliver an end-to-end solution that spans from low-level physical device configuration up to high-level algorithmic path computation and graphical user interaction.
\section{Organization of the Thesis}
The remainder of this thesis is logically structured to guide the reader through the theoretical background, the proposed algorithmic solutions, the system implementation, and the final evaluation. The document is organized as follows:
\begin{itemize}
    \item \textbf{Chapter 2:} Background and Literature Review provides the foundational concepts necessary to understand the research. It discusses the principles of Software-Defined Networking (SDN) in optical networks, introduces the TeraFlowSDN architecture, and details Optical Transport Networks (OTN) including WDM technology and ITU standards. The chapter concludes with a review of existing Routing and Spectrum Allocation (RSA) algorithms, highlighting the specific problems and limitations associated with parallel links.
    \item \textbf{Chapter 3:} Problem Definition and Proposed Algorithm formally establishes the research problem and introduces the mathematical formulation of the challenge. It details the graph modeling techniques used to represent topologies with parallel links and breaks down the proposed RSA algorithm, covering bandwidth computation, path discovery, bitmap-based spectrum allocation, and slot acquisition, followed by a thorough complexity analysis.
    \item \textbf{Chapter 4:} System Design and Implementation delves into the practical engineering and software architecture of the proposed solution. It outlines the modifications made to TeraFlowSDN—including device service enhancements for OpenConfig and OpenROADM, database schema extensions, and the RSA Computation Module. Furthermore, it details the development of the new Parallel Optical Controller microservice and the corresponding WebUI enhancements for spectrum visualization.
    \item \textbf{Chapter 5:} Testbed Setup and Deployment outlines the experimental environment used to validate the proposed system. It describes the physical and virtual infrastructure, the configuration of emulated transponders and ROADM devices, the setup of network topologies, and the deployment procedure utilizing MicroK8s.
    \item \textbf{Chapter 6:} Results and Evaluation presents a comprehensive analysis of the system's performance. It covers functional validation, computational performance, and spectrum utilization efficiency. The chapter also provides a throughput analysis, a comparative study against baseline approaches, and a statistical discussion of the results.
    \item \textbf{Chapter 7:} Conclusion and Future Work summarizes the primary contributions of the thesis, discusses the current limitations of the implemented system, and suggests potential directions for future research and development.
\end{itemize}